% Capitolo standalone: Reti Neurali
% (Non include preambolo né \documentclass; pensato per essere \input{} o \include{} nel main.)

\chapter{Reti Neurali}
\label{chap:neural-networks}

\section{Obiettivo e contesto}
L'obiettivo delle reti neurali è modellare relazioni potenzialmente non lineari tra gli indicatori di salute e la variabile target binaria \texttt{Diabetes\_binary}.
Il dataset presenta un forte sbilanciamento di classe (classe negativa maggioritaria), tipico di problemi di screening: di conseguenza, oltre all'accuratezza complessiva, diventano centrali metriche come \emph{recall} della classe positiva e la capacità di separazione misurata tramite ROC--AUC.

\section{Preprocessing e preparazione dei dati}
\subsection{Standardizzazione delle feature (\texttt{StandardScaler})}
Tutte le feature numeriche vengono standardizzate (media 0, varianza 1).
\textbf{Giustificazione:}
\begin{itemize}
  \item le reti neurali ottimizzate con metodi di discesa del gradiente sono sensibili alla scala delle feature: feature su scale molto diverse possono dominare il gradiente e rallentare/instabilizzare la convergenza;
  \item la standardizzazione rende più omogeneo il paesaggio della loss, facilitando l'ottimizzazione con Adam;
  \item è un prerequisito naturale per confronti con PCA, che è intrinsecamente sensibile alla scala.
\end{itemize}

\subsection{Suddivisione train/test}
Si adotta una suddivisione train/test 70/30 (\texttt{test\_size=0.3}) con \texttt{random\_state=42}.
\textbf{Giustificazione:}
\begin{itemize}
  \item 70/30 è un compromesso comune: mantiene un training set ampio per stimare i parametri e un test set abbastanza grande per una stima robusta delle metriche;
  \item fissare il seme casuale rende l'esperimento riproducibile.
\end{itemize}

\section{Scelte di valutazione e parametri operativi}
\subsection{Soglia di classificazione (\texttt{threshold=0.5})}
Le reti producono una probabilità $\hat{p}(y=1\mid x)$ tramite attivazione sigmoide.
La conversione in label predetta avviene con soglia $t=0.5$.
\textbf{Giustificazione:}
\begin{itemize}
  \item $t=0.5$ è la scelta standard quando la probabilità è calibrata e i costi di errore sono simmetrici;
  \item in contesto clinico i costi non sono simmetrici (i falsi negativi sono più gravi): mantenere $t=0.5$ come baseline permette un confronto equo tra configurazioni; eventualmente $t$ può essere ottimizzata per massimizzare recall o F$_\beta$ in una fase successiva.
\end{itemize}

\subsection{Metriche e diagnostica}
Per ogni modello vengono analizzati:
\begin{itemize}
  \item matrice di confusione e \texttt{classification\_report} (precision/recall/F1), per capire il compromesso tra falsi positivi e falsi negativi;
  \item curva ROC e area sotto la curva (ROC--AUC), per valutare la separabilità delle classi indipendentemente dalla soglia;
  \item curve di training/validation (loss e accuracy), per diagnosticare overfitting/underfitting.
\end{itemize}
\textbf{Giustificazione:} in presenza di sbilanciamento, l'accuratezza da sola può essere fuorviante (un modello che predice sempre la classe maggioritaria può avere accuracy elevata ma recall pessimo per la classe positiva).

\section{Architettura di base (No PCA, dataset originale)}
\subsection{Struttura}
La rete neurale baseline è una MLP (Multi-Layer Perceptron) fully-connected:
\begin{center}
\texttt{Input (21)} $\rightarrow$ \texttt{Dense(128, ReLU)} $\rightarrow$ \texttt{Dropout(0.3)} \\
$\rightarrow$ \texttt{Dense(64, ReLU)} $\rightarrow$ \texttt{Dropout(0.3)} \\
$\rightarrow$ \texttt{Dense(32, ReLU)} $\rightarrow$ \texttt{Dropout(0.3)} \\
$\rightarrow$ \texttt{Dense(1, Sigmoid)}
\end{center}

\subsection{Numero di neuroni: 128 $\rightarrow$ 64 $\rightarrow$ 32}
\textbf{Giustificazione:}
\begin{itemize}
  \item una struttura decrescente implementa una compressione progressiva: i primi layer apprendono combinazioni non lineari delle feature, mentre gli ultimi sintetizzano tali rappresentazioni verso una decisione binaria;
  \item 128 neuroni nel primo layer sono sufficienti per modellare interazioni tra 21 feature senza introdurre complessità eccessiva;
  \item la riduzione graduale (128$\to$64$\to$32) evita colli di bottiglia troppo bruschi (che potrebbero perdere informazione) e limita il numero di parametri, riducendo il rischio di overfitting.
\end{itemize}

\subsection{Attivazione ReLU nei layer nascosti}
\textbf{Giustificazione:}
\begin{itemize}
  \item ReLU è lo standard per MLP moderne grazie a efficienza computazionale e gradiente non saturante per input positivi;
  \item riduce il problema del \emph{vanishing gradient} rispetto a sigmoide/tanh nei layer nascosti;
  \item in assenza di evidenze di ``dead ReLU'' rilevanti su questo task, non è stato necessario introdurre varianti (LeakyReLU/ELU), mantenendo il modello più semplice e interpretabile.
\end{itemize}

\subsection{Dropout: $p=0.3$ dopo ogni layer}
\textbf{Giustificazione:}
\begin{itemize}
  \item il dropout introduce regolarizzazione stocastica, forzando la rete a non dipendere eccessivamente da subset specifici di neuroni;
  \item con dataset sbilanciato, la rete può overfittare pattern fortemente correlati alla classe maggioritaria; il dropout mitiga questo rischio;
  \item $p=0.3$ è un valore moderato: abbastanza alto da regolarizzare, ma non così aggressivo da impedire l'apprendimento.
\end{itemize}

\subsection{Output sigmoide}
\textbf{Giustificazione:}
\begin{itemize}
  \item un singolo neurone con sigmoide produce $\hat{p}(y=1\mid x)\in(0,1)$, interpretabile come probabilità;
  \item è la scelta naturale per classificazione binaria e si combina direttamente con la loss di cross-entropia binaria.
\end{itemize}

\subsection{Loss: Binary Cross-Entropy}
Con $y\in\{0,1\}$ e probabilità predetta $\hat{p}$:
\begin{equation}
\mathcal{L}_{BCE}(y,\hat{p}) = -\big(y\log \hat{p} + (1-y)\log(1-\hat{p})\big).
\end{equation}
\textbf{Giustificazione:}
\begin{itemize}
  \item è la loss standard per massima verosimiglianza in classificazione binaria (Bernoulli);
  \item penalizza in modo consistente predizioni molto confidenti ma errate, utile in contesti a rischio clinico.
\end{itemize}

\subsection{Ottimizzatore: Adam}
\textbf{Giustificazione:}
\begin{itemize}
  \item Adam usa stime adattive di primo e secondo momento del gradiente, riducendo la necessità di tuning manuale del learning rate;
  \item è robusto su problemi non lineari con feature standardizzate e dimensione campionaria medio-grande;
  \item rispetto a SGD ``puro'', in genere converge più rapidamente con meno iperparametri da tarare.
\end{itemize}

\subsection{Epoche: 150 \;\;\; Batch size: 256 \;\;\; Validation split: 0.1}
\textbf{Giustificazione:}
\begin{itemize}
  \item 150 epoche consentono di raggiungere una convergenza stabile monitorando le curve di validazione: il valore è scelto per lasciare margine oltre la fase di assestamento senza richiedere procedure aggiuntive (es. early stopping);
  \item batch size 256 è un compromesso tra stabilità del gradiente (batch più grandi) ed effetto regolarizzante del rumore (batch più piccoli), mantenendo efficienza computazionale;
  \item validation split 0.1 fornisce un set interno sufficiente per diagnosticare overfitting senza sottrarre troppi dati al training.
\end{itemize}

\section{Gestione dello sbilanciamento}
Lo sbilanciamento di classe richiede tecniche specifiche per evitare che la rete converga verso una soluzione ``comoda'' che predice prevalentemente la classe maggioritaria.
Sono stati sperimentati tre approcci: undersampling, class weights e focal loss.

\subsection{Undersampling + class balancing}
\subsubsection{Costruzione del dataset bilanciato}
Si seleziona casualmente dalla classe maggioritaria un numero di esempi pari a quelli della classe minoritaria, ottenendo una distribuzione 50/50.
\textbf{Parametri e scelte:}
\begin{itemize}
  \item campionamento \emph{senza rimpiazzo}: evita duplicazioni della classe maggioritaria e mantiene diversità;
  \item shuffle degli indici: evita bias dovuto all'ordine dei campioni;
  \item nuovo split train/test 70/30 con \texttt{stratify} sul dataset bilanciato: preserva il bilanciamento anche nelle partizioni.
\end{itemize}
\textbf{Giustificazione:} il bilanciamento ``a monte'' rende il segnale della classe positiva più presente durante il training, tipicamente aumentando il recall della classe positiva.
Lo svantaggio principale è la perdita di informazione (molti esempi della classe maggioritaria non vengono usati).

\subsubsection{Architettura ridotta sul dataset bilanciato}
Sul dataset ridotto si usa una rete meno capiente:
\begin{center}
\texttt{Input (21)} $\rightarrow$ \texttt{Dense(64, ReLU)} $\rightarrow$ \texttt{Dropout(0.5)} \\
$\rightarrow$ \texttt{Dense(32, ReLU)} $\rightarrow$ \texttt{Dropout(0.3)} \\
$\rightarrow$ \texttt{Dense(16, ReLU)} $\rightarrow$ \texttt{Dropout(0.3)} \\
$\rightarrow$ \texttt{Dense(1, Sigmoid)}
\end{center}
\textbf{Giustificazione dei parametri:}
\begin{itemize}
  \item neuroni dimezzati (64--32--16): dataset più piccolo $\Rightarrow$ maggior rischio di overfitting $\Rightarrow$ capacità ridotta;
  \item dropout 0.5 nel primo layer: regolarizzazione più forte dove il numero di connessioni è più alto;
  \item batch size 128 (ridotto) ed epoche 100 (ridotte): su dataset più piccolo si preferiscono aggiornamenti più frequenti e un training meno lungo per limitare overfitting.
\end{itemize}

\subsection{Class Weights sul dataset originale}
\subsubsection{Idea}
Si mantiene l'intero dataset ma si modifica la loss dando un peso maggiore agli errori sulla classe minoritaria (positiva).
In pratica, durante il training, ogni esempio contribuisce alla loss con un fattore dipendente dalla sua classe.
\textbf{Giustificazione:}
\begin{itemize}
  \item non si perdono campioni (a differenza dell'undersampling);
  \item si indirizza l'ottimizzazione verso un compromesso che penalizza maggiormente i falsi negativi;
  \item i pesi ``balanced'' dipendono dalla frequenza delle classi nel training set, evitando scelte arbitrarie.
\end{itemize}

\subsubsection{Architettura}
L'architettura resta quella della baseline (128--64--32, dropout 0.3), per isolare l'effetto dei class weights.
\textbf{Giustificazione:} mantenendo architettura e principali iperparametri invariati, le differenze in performance si attribuiscono più chiaramente alla sola strategia di pesatura delle classi.

\subsection{Focal Loss sul dataset originale}
\subsubsection{Motivazione}
La focal loss è progettata per problemi sbilanciati: riduce l'impatto degli esempi ``facili'' (correttamente classificati con alta confidenza) e concentra l'apprendimento sugli esempi difficili.

\subsubsection{Definizione (caso binario)}
Indicando con $p_t$ la probabilità del ``vero'' label:
\begin{equation}
 p_t = \begin{cases}
 \hat{p} & \text{se } y=1,\\
 1-\hat{p} & \text{se } y=0,
 \end{cases}
\end{equation}
la focal loss è:
\begin{equation}
\mathcal{L}_{FL}(y,\hat{p}) = -\alpha_t (1-p_t)^\gamma \log(p_t),
\end{equation}
con:
\begin{equation}
\alpha_t = \begin{cases}
 \alpha & \text{se } y=1,\\
 1-\alpha & \text{se } y=0.
 \end{cases}
\end{equation}
\textbf{Parametri scelti:}
\begin{itemize}
  \item $\gamma=2$: valore standard in letteratura, aumenta la focalizzazione sugli esempi difficili senza rendere l'ottimizzazione eccessivamente instabile;
  \item $\alpha$ calcolato dai dati come frequenza della classe negativa nel training set, $\alpha = \frac{n_0}{n_0+n_1}$.
\end{itemize}
\textbf{Giustificazione di $\alpha$:}
\begin{itemize}
  \item dato che la classe positiva (1) è minoritaria, assegnare $\alpha$ grande alla classe positiva (tramite $\alpha_t$) equivale a pesare di più gli errori sui positivi;
  \item calcolare $\alpha$ dalla distribuzione osservata evita un tuning manuale arbitrario e rende il setup replicabile.
\end{itemize}
\textbf{Nota operativa:} la focal loss introduce un costo computazionale maggiore rispetto a BCE/class weights e può richiedere più attenzione nel tuning; per questo motivo viene considerata principalmente come baseline di confronto.

\section{Confronto PCA vs No PCA nelle reti neurali}
Poiché la PCA non produce una riduzione drastica (serve un numero elevato di componenti per raggiungere l'80\% di varianza), viene usata come pipeline alternativa per confronto.

\subsection{Rete con input PCA e class weights}
Con input ridotto a $n_{comp}$ componenti, si adotta una rete più piccola:
\begin{center}
\texttt{Input ($n_{comp}$)} $\rightarrow$ \texttt{Dense(32, ReLU)} $\rightarrow$ \texttt{Dropout(0.3)} \\
$\rightarrow$ \texttt{Dense(16, ReLU)} $\rightarrow$ \texttt{Dropout(0.3)} \\
$\rightarrow$ \texttt{Dense(8, ReLU)} $\rightarrow$ \texttt{Dropout(0.3)} \\
$\rightarrow$ \texttt{Dense(1, Sigmoid)}
\end{center}
\textbf{Giustificazione:}
\begin{itemize}
  \item dimensionalità di input minore $\Rightarrow$ minore necessità di capacità del modello;
  \item architettura ridotta limita il rischio di overfitting e mantiene tempi di training comparabili;
  \item class weights restano utili per gestire lo sbilanciamento anche nello spazio PCA.
\end{itemize}

\subsection{Focal loss su input PCA}
Analogamente, la focal loss può essere applicata con gli stessi parametri $\gamma$ e $\alpha$ calcolato dalla distribuzione nel training set PCA.
\textbf{Giustificazione:} mantiene coerenza sperimentale e permette di valutare se la trasformazione PCA renda più/meno utili tecniche focalizzate sugli esempi difficili.

\section{Sintesi delle scelte}
\begin{itemize}
  \item \textbf{MLP fully-connected}: adatta a feature tabulari (indicatori di salute) e baseline solida.
  \item \textbf{ReLU + Dropout}: buon compromesso tra capacità e regolarizzazione.
  \item \textbf{Adam + BCE}: setup standard, stabile e replicabile.
  \item \textbf{Sbilanciamento}: confrontati undersampling (semplice ma informazione persa), class weights (usa tutti i dati), focal loss (focalizza sugli esempi difficili, più costosa).
  \item \textbf{PCA}: usata solo per confronto, con architetture coerentemente ridotte e stesse logiche di gestione dello sbilanciamento.
\end{itemize}
